\subsection{非零CGG族の射影圏の完全な離散分類}

\begin{theorem}[本族内部での射影圏の同値分類]
\label{thm:CGG-categorical-classification-v35}
$m=2j+r\ge3$、$j,r\ge0$、$s\ne0$ とし、
\[
 \mathcal Q_{m,j}(\alpha,s)
 :=\operatorname{qgr}\mathcal S_{j,m-2j}(\alpha,s)
\]
と書く。同様の条件を満たす二組のパラメータに対し、次は同値である。
\begin{enumerate}
\item $m=m'$ かつ $j=j'$。
\item $\mathcal Q_{m,j}(\alpha,s)$ と
 $\mathcal Q_{m',j'}(\alpha',s')$ は $\mathbb C$ 線形アーベル圏として同値である。
\item その有界導来圏は $\mathbb C$ 線形三角圏として同値である。
\end{enumerate}
従って固定した重み $m$ で得られる相異なる射影圏の個数は
$\lfloor m/2\rfloor+1$ である。
\end{theorem}

\begin{proof}
同じ $(m,j)$ では、全ての $\alpha$ がZhang twistで結ばれ、
全ての非零 $s$ が次数付き同型で結ばれるので、(1)から(2)が従う。
(2)から(3)は同値を複体に適用すればよい。

(3)を仮定する。標準full exceptional collectionの長さは $m+3$ であり、
その類はGrothendieck群の基底をなすので
\[
 K_0\bigl(D^b(\mathcal Q_{m,j}(\alpha,s))\bigr)
 \simeq\mathbb Z^{m+3}.
\]
従って $m=m'$ である。
通常の三角同値が存在すれば、そのtilting対象の像とRickardの定理により
二つの有限Beilinson代数の間にdg Morita同値も存在する。

両方の $r=m-2j,r'=m-2j'$ が正なら、
定理\ref{thm:fixed-quantum-hull-v35}によるminiversal環の長さは
それぞれ $j+1,j'+1$ であり、導来不変性から $j=j'$ を得る。
一方のみ $r=0$ なら、定理\ref{thm:endpoint-hull-v36}により
そのhullは次元 $m+2$ の正則環である一方、他方のhullはArtin環である。
これはhullの導来不変性に矛盾する。
両方が零なら $j=j'=m/2$ である。これで(3)から(1)が従う。
\end{proof}

この分類の対象は、明示した接線共鳴CGG族とその $r=0$ 端点である。
全ての非可換重み付き射影空間の分類を主張するものではない。
重み三の四対象VK部分圏は、第\ref{sec:vk-intrinsic-classification-v36}節で別に分類する。

\begin{corollary}[可換点を含めた族]固定した $m$ において $s=0$ も許すと、
本族の射影圏は正確に $\lfloor m/2\rfloor+2$ 個の導来同値類を持つ。
\end{corollary}
\begin{proof}$s=0$ では全て同じ可換射影スタックの圏になる。
その $HH^1$ の次元は $N_m+9$、非零量子化では一であるため、
可換圏は全ての非零族から区別される。\end{proof}
